% !TeX root = ../Tex/main.tex

\begin{center}
    \subsection*{Abstract}
\end{center}

Standard game-theoretic accounts often rationalize departures from fully rational play by appealing to exogenous "shocks" or moments of surprise. Yet many strategic encounters suggest that what matters is not merely an external disturbance but an endogenous shift that reshapes perception, signalling, and restraint. This paper develops an illustrative framework using two elephant-human encounters reported by Safina (2015) as motivating cases. In one episode, an elephant charges a tourist yet brakes before contact; in another an elephant injures a shepherd and subsequently guards him until help arrives. Canonical conflict models can capture incentives and credible threats in both interactions, but they struggle to accommodate the joint pattern "attack, then care." I treat the two episodes as potential outcomes of a common underlying interaction, and use a causal directed acyclic graph to make explicit the missing variables that standard models implicitly omit. I argue that surprise operates as an informational shock that can trigger aggression, while empathy functions as an affective regulator that modulates signal intensity and can generate post-conflict assistance. To formalise this intuition, I sketch an adaptation of Groseclose and Snyder's (1996) vote-buying game, in which competing emotions allocate limited intensity across competing "impulses" to determine behaviour. The resulting perspective offers a tractable route to richer models of bounded rationality and yields testable implications about when shocks produce escalation versus restraint.

% \lipsum[1] % genera il primo paragrafo di testo Lorem Ipsum